\label{sec:conclusion}

The redistricting criterion ``preserve communities of interest'' is, at present,
a phrase that courts accept and litigants invoke but that no judge can quantitatively
enforce and no algorithm can natively satisfy.
This paper has proposed a mathematical framework that changes that.
By representing community character as a vector of tract-level signals, by measuring
community similarity as cosine similarity of those vectors, and by encoding similarity
as an METIS edge weight modifier, the framework translates a qualitative legal
criterion into a deterministic, verifiable, partisan-neutral computation.

The framework satisfies the key legal requirements for defensible redistricting criteria.
It is neutral under \textit{Shaw v.\ Reno} and \textit{Miller v.\ Johnson}: no
racial, partisan, or incumbency data enters the character vector.
It is consistent with the communities of interest doctrine as articulated in
\textit{Reynolds v.\ Sims} and as codified in California, Colorado, Arizona, and
Washington state law.
It is quantitative: the formula $w(u,v) = w_{\mathrm{base}} \times \mathrm{sim}(\mathbf{c}_u, \mathbf{c}_v)$
is fully disclosed, replicable by any party with access to the same public data, and
immune to the charge that community identifications were made ad hoc for political advantage.

\subsection{Expert Witness Role}

The framework is designed for use by expert witnesses in redistricting litigation.
A cartographer or political scientist serving as an expert witness can testify:
\begin{itemize}
  \item The district boundaries were placed by an algorithm that minimizes cuts
    through edges connecting tracts of high community similarity.
  \item Similarity was computed from [NLCD / LODES / TIGER / EIA 861] data, which is
    publicly available and was processed by a deterministic formula.
  \item The formula does not incorporate, and is not aware of, the racial composition,
    partisan registration, or prior voting history of any tract.
  \item The communities preserved by the algorithm correspond to communities
    recognized in prior redistricting litigation and commission proceedings as
    legitimate communities of interest: school district co-members, commuting-shed
    partners, economic sector colleagues.
\end{itemize}
This template, developed in the context of M-track analysis, provides a standard
framework for the expert declaration that would accompany any plan submitted
to a court or commission.

\subsection{Relationship to M.1--M.8}

M.0 is the theoretical reference; M.1--M.7 are the empirical papers.
Each of M.1--M.7 inherits the framework defined here---character vectors, cosine
similarity, edge weight modifier, alpha-blending---and specializes it to one
data source.
Each paper validates the specialization by running \textsc{bisect} on a selection
of states and measuring the effect of the character weights on plan outcomes:
compactness, school district integrity, administrative zone preservation, and
proportionality gap.
M.8 will compose all seven signals into a single composite index and provide
an end-to-end expert witness template.

\subsection{Open Questions}

Three open questions are not addressed in this paper and are left for future work:
\begin{enumerate}
  \item \textbf{Optimal alpha}: The alpha-blending parameter $\alpha$ is currently
    a fixed configuration choice.
    A data-driven method for selecting $\alpha$ to maximize community preservation
    subject to compactness constraints would be valuable.
  \item \textbf{Signal orthogonality}: The M.1--M.7 signals are not independent;
    for example, economic character and commuting shed are correlated.
    A principal component analysis of the joint character matrix could identify
    a smaller set of orthogonal community dimensions.
  \item \textbf{Temporal stability}: Community character changes over time.
    Whether a weight regime calibrated on the 2010 census remains appropriate for
    the 2020 redistricting cycle is an empirical question not addressed here.
\end{enumerate}

The community character framework proposed in this paper provides a rigorous,
legally defensible, and computationally tractable answer to the question that
has long vexed redistricting reform: what, precisely, is a community of interest,
and how do you know whether you have preserved one?
